\section{Multilinear estimates and Fourier lifts}
\label{sec:lifts}

The tuple laws of Section~\ref{sec:tuples} provide two ways to control
multilinear averages: reflection positivity compares different inputs
with a single diagonal quantity, while conditional mixing suppresses
large Fourier denominators.  We develop both estimates, then prove the
moment bounds needed to apply them to bounded functions on integer
intervals.  The signed-form estimates adapt~\cite[Proposition~4.1 and
Lemmas~4.2--4.3]{squarepower}.  We prove the needed estimates for the
present laws, with changes of auxiliary parameter replacing the dilation
invariance used in the square setting.

Retain the index set \(V=\Pi^w\), where \(\Pi\) consists of the
permutation words on \([L]\), and put \(d=|V|\); in particular, \(d\) is even.
Fix a positive integer
\(l\) and a finite set \(\mathcal P\) of primes at least \(P\).  Set
\[
 X_{\mathcal P}=\prod_{p\in\mathcal P}\F_p,
 \qquad
 \mu_{l,\mathcal P}=\bigotimes_{p\in\mathcal P}\mu_{l,p}.
\]
We collect the prime labels in slot \(v\) into a single variable
\(z_v\in X_{\mathcal P}\), so \(\mu_{l,\mathcal P}\) is a probability
law on \(X_{\mathcal P}^V\).  Each single-slot marginal is uniform.
Norms of functions on \(X_{\mathcal P}\) will always use uniform
probability measure unless another measure is specified.
For real functions on this space define
\[
 \mathcal Z_{l,\mathcal P}((g_v)_{v\in V})
 =\int\prod_{v\in V}g_v(z_v)\,d\mu_{l,\mathcal P},
 \qquad
 Z_{l,\mathcal P}(g)=\mathcal Z_{l,\mathcal P}((g)_{v\in V}).
\]
The empty prime set is allowed: its residue space is a singleton and
\(Z_{l,\varnothing}(g)=g^d\).

We also fix the Fourier notation on these spaces.  Every frequency has a
unique additive decomposition
\[
 \xi=\sum_{p\in\mathcal P}\frac{a_p}{p}\pmod 1,
 \qquad 0\le a_p<p,
\]
and its character is
\(e(z\xi)=\prod_{p\in\mathcal P}e(a_pz_p/p)\).
Equivalently, one may use any integer representative of \(z\) supplied
by the Chinese remainder theorem.  The \emph{exact support} of this
character is \(\{p:a_p\ne0\}\); its reduced denominator is the product
of those primes.  For \(B\subseteq\mathcal P\) write
\(q_B=\prod_{p\in B}p\), with \(q_\varnothing=1\).

\subsection{Signed H\"older inequality and seminorm rules}

The finite reflection argument below follows the chessboard method of
Fr\"ohlich, Israel, Lieb and Simon~\cite[Theorem~4.1]{FILS1978};
related reflection arguments for graph norms appear in Conlon and
Lee~\cite[Section~3]{ConlonLee2017}.
The standard passage from a multilinear H\"older inequality to a seminorm
is also part of Hatami's criterion for graph norms~\cite[Theorem~2.8 of
the arXiv version]{Hatami2010}; we prove both steps directly for the
present tuple laws.

\begin{proposition}[Signed H\"older inequality]
\label{lift:signed-holder}
For every positive integer \(l\), finite set \(\mathcal P\) of primes
at least \(P\), and real functions \(g,g_v:X_{\mathcal P}\to\R\),
\begin{equation}
 Z_{l,\mathcal P}(g)\ge0,
 \qquad
 \left|\mathcal Z_{l,\mathcal P}((g_v)_{v\in V})\right|
 \le\prod_{v\in V} Z_{l,\mathcal P}(g_v)^{1/d}.
 \label{eq:10}
\end{equation}
\end{proposition}

\begin{proof}
Across any reflection cut of \(V\), the matrix of the product law is
the tensor product of the corresponding prime-law matrices.  It is
therefore positive semidefinite.  To explain its consequence for signed
inputs, fix a cut \(V=V_+\sqcup V_-\) and its reflection \(\theta\).
Define maps \(\phi_\pm:V\to V_\pm\) that fix the selected half and
reflect the other half into it.  An assignment of functions is a map
\(a:v\mapsto g_v\); write its integral as \(\Lambda(a)\).
The two real functions on the plus-half labels that enter the cut
bilinear form are
\[
 F(x)=\prod_{v\in V_+}a(v)(x_v),
 \qquad
 G(x)=\prod_{v\in V_+}a(\theta v)(x_v).
\]
Its mixed value is \(\Lambda(a)\), and its two quadratic values are
\(\Lambda(a\circ\phi_+)\) and \(\Lambda(a\circ\phi_-)\).
Cauchy--Schwarz for the positive semidefinite matrix gives
\begin{equation}
 |\Lambda(a)|^2
 \le \Lambda(a\circ\phi_+)\Lambda(a\circ\phi_-),
 \qquad \Lambda(a\circ\phi_\pm)\ge0.
 \label{lift:fold-cs}
\end{equation}
In particular a constant assignment has nonnegative integral, proving
the first assertion of~\eqref{eq:10}.

We next give a finite sequence of fold maps that sends every vertex to
one vertex.  In one block write \(p_i\) for the position of symbol
\(i\) in its permutation word.  Folding onto the half where \(i\)
precedes \(i+1\) sorts the pair \((p_i,p_{i+1})\), leaving all other
entries unchanged.  Apply the comparisons \(i=1,\ldots,L-1\) in that
order, and repeat this sweep \(L-1\) times.  A sweep moves the largest
entry to the end.  Once the largest entries occupy their final positions,
the next sweep leaves them there and places the largest remaining entry
immediately before them.  Thus every position array becomes
\((1,\ldots,L)\).  Perform these sweeps in each block.  If their fold
maps in chronological order are \(\phi_1,\ldots,\phi_m\), then
\[
 \phi_m\circ\cdots\circ\phi_1(v)=v_*
 \quad\text{for every }v\in V,
\]
where \(v_*\) has the identity word in every block.

Now fix a finite nonempty list of real functions and let \(M\) be the
maximum absolute integral over all assignments from that list, allowing
repetition.  If \(M>0\), take a maximizing assignment.  Each of its
folded integrals is nonnegative and at most \(M\), while their product
is at least \(M^2\) by~\eqref{lift:fold-cs}.  Both therefore equal
\(M\).  Hence either fold preserves maximality.  Apply the pullbacks
in reverse order, first \(\phi_m\), then \(\phi_{m-1}\), and so on.
The final assignment is
\(a\circ\phi_m\circ\cdots\circ\phi_1\), which is constant.  It
follows that \(M\) is a diagonal value.  Conversely, constant
assignments are included in the maximum.  We have proved that the
maximum absolute mixed integral equals the largest diagonal value in
the list; this is also true when \(M=0\).

For \(\epsilon>0\), apply this conclusion to the list
\[
 \frac{g_v}{(Z_{l,\mathcal P}(g_v)+\epsilon)^{1/d}},
 \qquad v\in V.
\]
Each diagonal value in this list is at most one.  Multilinearity gives
\( |\mathcal Z_{l,\mathcal P}((g_v)_v)|
 \le\prod_v(Z_{l,\mathcal P}(g_v)+\epsilon)^{1/d}\).
Let \(\epsilon\downarrow0\).  This proves the second assertion,
including cases where some diagonal values vanish.
\end{proof}

Define
\[
 \|g\|_{l,\mathcal P}=Z_{l,\mathcal P}(g)^{1/d}
 \quad\text{for real }g.
\]
For \(B\subseteq\mathcal P\), let \(\E_B g\) mean uniform averaging
over the coordinates in \(B\), with all other coordinates held fixed.
The resulting function may be viewed on either \(X_{\mathcal P}\) or
\(X_{\mathcal P\setminus B}\).

\begin{lemma}[Seminorm and changes of coordinates]
\label{lift:seminorm}
The map \(g\mapsto\|g\|_{l,\mathcal P}\) is a seminorm on the real
functions on \(X_{\mathcal P}\).  It satisfies
\[
 |\E g|\le\|g\|_{l,\mathcal P},
 \qquad
 \|\E_Bg\|_{l,\mathcal P\setminus B}
 =\|\E_Bg\|_{l,\mathcal P}
 \le\|g\|_{l,\mathcal P}.
\]
Translations of the input are isometries.  Moreover, if \(D\) is a
positive integer coprime to every prime of \(\mathcal P\), then for
any \(a'\in X_{\mathcal P}\),
\begin{equation}
 \left\|g\bigl(\lambda(D)^{-1}(\,\cdot-a')\bigr)
       \right\|_{l,\mathcal P}
 =\|g\|_{lD,\mathcal P}.
 \label{lift:transport}
\end{equation}
Here the affine change is interpreted separately at each prime.
\end{lemma}

\begin{proof}
Nonnegativity follows from Proposition~\ref{lift:signed-holder}.
Since \(d\) is even, homogeneity of \(Z\) gives absolute homogeneity
of its \(d\)-th root.  Expanding \(Z(g+h)\) by its slots and applying
\eqref{eq:10} termwise gives
\[
 Z_{l,\mathcal P}(g+h)
 \le\sum_{S\subseteq V}\|g\|_{l,\mathcal P}^{|S|}
                         \|h\|_{l,\mathcal P}^{d-|S|}
 =\bigl(\|g\|_{l,\mathcal P}+\|h\|_{l,\mathcal P}\bigr)^d.
\]
Taking roots proves the triangle inequality.
Putting \(g\) in one slot and 1 in all others in~\eqref{eq:10}
proves mean domination, since the single-slot marginal is uniform and
\(Z(1)=1\).

The simultaneous translation invariance of the tuple laws proves that
translations are isometries.  Let \(G_B\subseteq X_{\mathcal P}\)
be the additive subgroup of vectors supported on \(B\).  Then
\[
 (\E_Bg)(z)=\frac1{|G_B|}\sum_{a\in G_B}g(z+a).
\]
The triangle inequality gives the asserted contraction.  Since
\(\E_Bg\) is independent of \(B\), the probability laws at those
primes integrate out, proving the equality of the two seminorms.
Finally, \(\lambda(D)\) is a unit at each remaining prime, and
\eqref{eq:7} pushes the product law at \(l\) exactly to that at
\(lD\).  This proves~\eqref{lift:transport}.
\end{proof}

\subsection{Decay at large Fourier denominators}

Conditional mixing gives a different estimate, valid also after
removing the exceptional part of each prime law.  Its proof requires
orthogonality between different exact supports; applying the triangle
inequality to all Fourier coefficients would lose this decay.

\begin{lemma}[Large-denominator estimate]
\label{lift:large-denominator}
Fix \(l\ge1\) and a finite prime set \(\mathcal P\) as above.  At
each prime choose either \(\rho_p=\mu_{l,p}\) or
\(\rho_p=\mu_{l,p}^{\circ}\).  Let \(g_v:X_{\mathcal P}\to\C\)
be arbitrary functions.  If, in one distinguished slot \(v\), every
Fourier coefficient at a denominator at most \(R\) vanishes, where
\(R\ge1\), then
\begin{equation}
 \left|\int\prod_{u\in V}g_u(z_u)\,
                  d\bigotimes_{p\in\mathcal P}\rho_p\right|
 \le R^{-\gamma}\|g_v\|_2
                     \prod_{u\ne v}\|g_u\|_{2(d-1)}.
 \label{eq:11}
\end{equation}
\end{lemma}

\begin{proof}
At one prime, let \(\mathfrak m_p\) denote the marginal law of the labels in all
slots other than \(v\), and let
\[
 T_p f=\E_{\rho_p}\bigl(f(z_v)\mid(z_u)_{u\ne v}\bigr).
\]
This is an operator from uniform \(L^2(\F_p)\) to \(L^2(\mathfrak m_p)\).
It preserves constants.  It sends uniform-mean-zero inputs to
\(\mathfrak m_p\)-mean-zero outputs, since \(\E_{\mathfrak m_p}T_pf=\E f\).
By~\eqref{eq:8}, its norm on the mean-zero subspace is at most
\(p^{-\gamma}\).

The product law makes the conditional operator for all primes equal
to \(T=\bigotimes_pT_p\): this identity holds first for products of
one-prime functions, by independence across primes, and then for every
function by linearity.  The exact-support \(B\) Fourier space has a
mean-zero factor at primes in \(B\) and a constant factor elsewhere.
On this space the operator norm is at most \(q_B^{-\gamma}\).
Furthermore, the images of two distinct exact-support spaces are
orthogonal in \(L^2(\bigotimes_p\mathfrak m_p)\).  At a prime in the symmetric
difference of the supports, one image factor is constant and the other
has mean zero.  Their inner product is zero, and tensoring preserves
this orthogonality.

Decompose \(g_v=\sum_{q_B>R}g_{v,B}\) by exact supports.  Both the
domain summands and their images are orthogonal, so
\[
 \|Tg_v\|_2^2
 =\sum_{q_B>R}\|Tg_{v,B}\|_2^2
 \le\sum_{q_B>R}q_B^{-2\gamma}\|g_{v,B}\|_2^2
 \le R^{-2\gamma}\|g_v\|_2^2.
\]
Condition the original integral on the other slots and apply
Cauchy--Schwarz.  The remaining factor is bounded by
\[
 \left\|\prod_{u\ne v}g_u(z_u)\right\|_2
 \le\prod_{u\ne v}
       \left(\E|g_u(z_u)|^{2(d-1)}\right)^{1/(2(d-1))}
 =\prod_{u\ne v}\|g_u\|_{2(d-1)},
\]
using ordinary H\"older and the uniform single-slot marginals.
This proves~\eqref{eq:11}.  If \(\mathcal P\) is empty, the frequency
hypothesis forces \(g_v=0\), and the conclusion is immediate.
\end{proof}

\subsection{Fourier lifts and uniform moments}

A bounded function on an integer interval gives Fourier coefficients at
rational frequencies.  We transfer a finite set of these coefficients
to \(X_{\mathcal P}\).  Such a transfer need not preserve pointwise
bounds, so we prove bounds for every fixed moment, following the
Fourier-lift argument of~\cite[Lemmas~5.1--5.2]{squarepower}.  We must also allow
arbitrary subfamilies of complete exact supports, because fixing some
residue coordinates changes which denominators remain.

Let \(I\) be a nonempty consecutive interval of integers and let
\(f:I\to\C\).  Put
\[
 \widehat f_I(\xi)=\frac1{|I|}\sum_{x\in I}f(x)e(-x\xi).
\]
For \(B\subseteq\mathcal P\), define its exact-support contribution
\[
 f_B(z)=\sum_{q(\xi)=q_B}\widehat f_I(\xi)e(z\xi),
\]
where the frequencies are characters of \(X_{\mathcal P}\).  For
\(Q\ge1\), the Fourier lift is
\[
 \mathcal L_{I,\mathcal P,Q}f(z)
 =\sum_{\substack{B\subseteq\mathcal P\\q_B\le Q}}f_B(z)
 =\sum_{\substack{\xi\text{ on }X_{\mathcal P}\\q(\xi)\le Q}}
                   \widehat f_I(\xi)e(z\xi).
\]
Each \(f_B\) depends only on the coordinates in \(B\), and has mean
zero in each of those coordinates.  If \(f\) is real, all these
functions are real because each summation set is closed under negation.
Retaining a support \(B\) always means retaining all its frequencies.

\begin{proposition}[Uniform lift moments]
\label{lift:moments}
Let \(\mathcal P\) be any finite set of primes, \(Q\ge1\), and let
\(I\) be a consecutive integer interval with \(|I|\ge10Q^6\).
For every integer \(r\ge1\), every \(\epsilon>0\), every function
\(f:I\to\C\), and every subfamily
\(\mathcal F\subseteq\{B\subseteq\mathcal P:q_B\le Q\}\),
\begin{equation}
 \left\|\sum_{B\in\mathcal F}f_B\right\|_{2r}
 \le C_{r,\epsilon}\|f\|_\infty Q^\epsilon.
 \label{eq:12}
\end{equation}
The constant is independent of the prime set, the interval and its
location, the function, and the subfamily.
\end{proposition}

\begin{proof}
The zero function is immediate; otherwise scale so that
\(\|f\|_\infty\le1\).  Write \(N_I=|I|\), and let \(u\) be the
largest cardinality of a support whose prime product is at most \(Q\).
We first control conditional sums of squares of the \(f_B\), then their
square-function moments.  Symmetrization will recover the moment of
their sum.

\smallskip
\noindent\emph{Conditional energy.}
For \(J\subseteq\mathcal P\) with \(q_J\le Q\), and any fixed
coordinate vector \(z_J\), we claim that
\begin{equation}
 \sum_{\substack{B\in\mathcal F\\B\supseteq J}}
       \E\bigl(|f_B|^2\mid z_J\bigr)\ll4^{|J|}.
 \label{eq:13}
\end{equation}
The conditional expectation here averages uniformly over all coordinates
outside \(J\).  If \(B\supseteq J\), the coefficient of an outside
character \(\eta\) of exact support \(B\setminus J\), after fixing
\(z_J\), is
\[
 c_J(\eta;z_J)
 =\frac1{N_I}\sum_{x\in I}f(x)e(-x\eta)
       \prod_{p\in J}\bigl(p\1_{\{x\equiv z_p\pmod p\}}-1\bigr).
\]
Indeed, the sum over the nonzero frequencies at \(p\) is
\(\sum_{a=1}^{p-1}e(a(z_p-x)/p)
=p\1_{\{x\equiv z_p\pmod p\}}-1\).
Different supports \(B\supseteq J\) have disjoint sets of outside
characters.  Conditional Parseval therefore identifies the left side of
\eqref{eq:13} with the squared Euclidean norm of this coefficient vector
on the union of those character sets.

Expand the product in \(c_J\) over subsets \(J'\subseteq J\).
Apart from a sign, the corresponding vector has coordinates
\[
 \frac{q_{J'}}{N_I}
 \sum_{\substack{x\in I\\x\equiv z_{J'}\pmod{q_{J'}}}}
                   f(x)e(-x\eta).
\]
The congruence denotes the class determined by the fixed coordinates in
\(J'\); for the empty subset it imposes no restriction.  Write its
progression as \(x=a+q_{J'}n\), with \(n\) in a consecutive interval
of length \(N'\).  Then
\[
 N'\ge0.9N_I/q_{J'},
 \qquad N'q_{J'}/N_I\le1.1.
\]
In these coordinates the frequencies are \(q_{J'}\eta\), up to
constant phases.  They remain distinct because their denominators are
coprime to \(q_J\).  Their denominators are at most \(Q\), so they
number at most \(Q^2\) and have pairwise circle distance at least
\(Q^{-2}\).

The Gram matrix of these characters in normalized \(L^2\) on the
progression has diagonal 1.  By geometric summation, every off-diagonal
entry has modulus at most \(Q^2/(2N')\).  Its operator norm is
therefore at most its maximum absolute row sum, which is bounded by
\[
 1+\frac{Q^4}{2N'}
 \le1+\frac{Q^5}{1.8N_I}
 \le1+\frac1{18Q}.
\]
The associated Fourier analysis map has norm squared equal to this
Gram-matrix norm.  Since \(\|f\|_\infty\le1\) and
\(N'q_{J'}/N_I\le1.1\), the displayed coefficient vector has
Euclidean norm bounded by an absolute constant.  The triangle inequality
over the \(2^{|J|}\) subsets \(J'\), followed by squaring, proves
\eqref{eq:13}.

\smallskip
\noindent\emph{Square-function moments.}
Set
\[
 S(z)=\left(\sum_{B\in\mathcal F}|f_B(z)|^2\right)^{1/2}.
\]
Expand \(\E S^{2r}\) and fix the supports of its first \(r-1\)
factors.  Their union \(W'\) has size at most \((r-1)u\), and their
product depends only on \(z_{W'}\).  For the final support \(B\),
integration outside \(W'\) leaves exactly
\(\E(|f_B|^2\mid z_J)\), where \(J=B\cap W'\).
Every intersection that occurs satisfies \(q_J\le q_B\le Q\).
For each fixed intersection \(J\), the sum over supports with
\(B\cap W'=J\) is bounded by the larger sum over all \(B\supseteq J\)
in~\eqref{eq:13}.  There are at most \(2^{(r-1)u}\) possible
intersections and each has at most \(u\) elements.  Thus
\[
 \E S^{2r}\le C\,2^{(r-1)u}4^u\E S^{2r-2}.
\]
The same argument with \(r=1\) uses \(J=W'=\varnothing\).
Iteration gives
\begin{equation}
 \|S\|_{2r}\le C_r^{1+u}.
 \label{lift:square-function}
\end{equation}
The overlaps among supports have now been controlled.  It remains to
bound their sum by this square function with a loss exponential only in
\(u\).

\smallskip
\noindent\emph{Recovering the sum.}
Fix a support size \(j\ge1\) and color every prime coordinate with one
of \(j\) colors.  Call a support \emph{transversal} when it has exactly
one prime of each color, and let \(\mathcal T\) be the transversal
supports in \(\mathcal F\) for this coloring.  Take independent uniform
pairs \(X_p^0,X_p^1\) at every prime.  The operator
\(\operatorname{swap}_p\) interchanges these two copies.  Define
\[
 D_B=\prod_{p\in B}(1-\operatorname{swap}_p)f_B(X^0),
 \qquad B\in\mathcal T.
\]
Averaging over all second copies leaves \(f_B(X^0)\), since every
other term in this expansion has mean zero in at least one copied
coordinate.  Jensen's inequality gives
\[
 \left\|\sum_{B\in\mathcal T}f_B\right\|_{2r}
 \le\left\|\sum_{B\in\mathcal T}D_B\right\|_{L^{2r}(X^0,X^1)}.
\]
Independent swaps preserve the joint law of the copies and multiply
\(D_B\) by the product of the corresponding signs.  Thus for independent
uniform signs \(\epsilon_p\), indexed by primes,
\[
 \left\|\sum_{B\in\mathcal T}D_B\right\|_{L^{2r}(X^0,X^1)}
 =\left\|\sum_{B\in\mathcal T}D_B\prod_{p\in B}\epsilon_p
       \right\|_{L^{2r}(X^0,X^1,\epsilon)}.
\]

For completeness, the scalar sign estimate
\[
 \left\|\sum_i a_i\epsilon_i\right\|_{2r}
 \le C_r\left(\sum_i|a_i|^2\right)^{1/2}
\]
follows by expanding the even moment.  Surviving terms have even
multiplicity at every index.  For nonnegative coefficients their sum is
at most
\[
 \frac{(2r)!}{2^r r!}\left(\sum_i a_i^2\right)^r,
\]
because \((2m)!\ge2^m m!\).  For complex coefficients, take absolute
values of the surviving terms in
\((\sum_i a_i\epsilon_i)^r(\sum_i\overline{a_i}\epsilon_i)^r\)
and use the nonnegative case.  Iterating this estimate over the color
classes gives, for fixed copies,
\[
 \left\|\sum_{B\in\mathcal T}D_B\prod_{p\in B}\epsilon_p
       \right\|_{L^{2r}(\epsilon)}
 \le C_r^j\left(\sum_{B\in\mathcal T}|D_B|^2\right)^{1/2}.
\]
At each step Minkowski is used in the form
\(\| (\sum_i|B_i|^2)^{1/2}\|_{2r}
\le(\sum_i\|B_i\|_{2r}^2)^{1/2}\).
The signs are indexed by individual primes; transversality ensures that
each monomial uses exactly one sign from each color class.

In expanding the swaps in \(D_B\), index its \(2^j\) terms by a choice
of copy for each color.  For any fixed such choice, made consistently
across all supports, the selected prime coordinates have the original
product law.  The triangle inequality in
\(L^{2r}(\ell^2(\mathcal T))\) therefore gives
\[
 \left\|\left(\sum_{B\in\mathcal T}|D_B|^2\right)^{1/2}\right\|_{2r}
 \le2^j\left\|\left(\sum_{B\in\mathcal T}|f_B|^2\right)^{1/2}
       \right\|_{2r}
 \le2^j\|S\|_{2r}.
\]
Consequently the transversal sum has norm at most
\((2C_r)^j\|S\|_{2r}\).

Now color independently and uniformly.  A fixed size-\(j\) support is
transversal with probability \(j!/j^j\ge e^{-j}\), where here \(e\)
is Euler's number; the inequality follows from
\(\log(j!)\ge\int_1^j\log x\,dx\).
Thus the full size-\(j\) sum is exactly the coloring average of its
transversal sum divided by this probability.  Minkowski yields
\[
 \left\|\sum_{\substack{B\in\mathcal F\\|B|=j}}f_B\right\|_{2r}
 \le(2eC_r)^j\|S\|_{2r}.
\]
Include the constant term of size zero and sum over \(j\le u\).
Together with~\eqref{lift:square-function}, this bounds the norm of
the full sum by \((C'_r)^{1+u}\).
Finally, the product of \(u\) distinct primes is at least \((u+1)!\),
so \((u+1)!\le Q\).  Hence \(u=o(\log Q)\) as \(Q\to\infty\),
uniformly in the prime set.  For every fixed \(r\) and \(\epsilon>0\),
\((C'_r)^{1+u}\le C_{r,\epsilon}Q^\epsilon\).
Restoring \(\|f\|_\infty\) proves~\eqref{eq:12} with all the stated
uniformities.
\end{proof}

\begin{lemma}[Specialization of residue coordinates]
\label{lift:specialization}
Let \(B\subseteq\mathcal P\), fix \(z_B\in X_B\), and let
\(G:X_{\mathcal P}\to\C\).  For \(1\le s<\infty\),
\[
 \|G(z_B,\cdot)\|_{L^s(X_{\mathcal P\setminus B})}
 \le q_B^{1/s}\|G\|_{L^s(X_{\mathcal P})}.
\]
In particular, a Fourier lift or a subfamily of its complete supports
satisfies~\eqref{eq:12} after specialization with the additional factor
\(q_B^{1/(2r)}\).  The same bound holds after retaining either the
remaining denominators at most a given \(H\ge1\), or those greater
than \(H\).
\end{lemma}

\begin{proof}
The specified fiber has uniform probability \(1/q_B\).  Its contribution
to \(\E|G|^s\) is
\(q_B^{-1}\E|G(z_B,\cdot)|^s\).  Dropping the other nonnegative
contributions proves the first assertion.  For the last assertion,
a character with original exact support \(C\) retains exact support
\(C\setminus B\) after specialization.  The two cuts therefore arise,
before specialization, from the complete-support subfamilies
\[
 \{C\in\mathcal F:q_{C\setminus B}\le H\},
 \qquad
 \{C\in\mathcal F:q_{C\setminus B}>H\}.
\]
Apply Proposition~\ref{lift:moments} to these subfamilies and then use
the fiber bound.  Combining coefficients can cancel frequencies but
cannot create any outside the indicated ranges.
\end{proof}

In particular, applying~\eqref{eq:11} with the same set of coordinates
fixed in every slot incurs total specialization factor
\[
 q_B^{1/2}\left(q_B^{1/(2(d-1))}\right)^{d-1}=q_B.
\]
Two specialized \(L^2\) norms incur the same factor.  All moment orders
needed here are fixed after \(h\) is fixed.  Finite products of their
subpower bounds remain subpower; none of these constants depends on
\(l\), the specialized residues, or the retained supports.

\subsection{Parameters and the density quantity}

We now choose the analytic parameters used in the kernel estimate and
the scale comparison.  The constants \(k,d,\gamma\) have already been
fixed.  Set
\begin{equation}
 \begin{aligned}
 \eta&=\frac1{100k},&
 \zeta&=\frac{\eta}{20k},&
 \nu&=\frac{\eta}{10\,2^k},\\
 \beta&=\frac{\nu}{100},&
 \tau&=\frac{\beta\gamma}{1000k(d+1)},&
 \sigma&=\frac{\gamma\tau}{4}.
 \end{aligned}
 \label{eq:14}
\end{equation}
The displayed parameters are functions only of the already fixed numbers
\(k,d,\gamma\).  When Condition~\eqref{local:unit-condition} holds and the
designated choices of Sections~\ref{sec:local} and~\ref{sec:tuples} are made,
these numbers, and hence every parameter in~\eqref{eq:14}, depend only on
\(k\).
Here \(\eta\) is distinct from the exceptional masses \(\eta_{l,p}\).
At an integer scale \(N\ge1\), the kernel proof will use \(\eta\) to
set the denominators and widths of neighborhoods of rational frequencies,
\(\nu\) for the exponential-sum saving \(N^{-\nu}\) outside them,
and \(\zeta\) for the allowed range \(l\le N^\zeta\).
The exponent \(\beta\) fixes the Fourier denominator cutoff \(N^\beta\);
\(\tau\) sets the subdivision into about \(N^\tau\) intervals and
the bound \(q_B\le N^\tau\) in the shorter-scale comparison, while
\(\sigma\) is the common error exponent in \(N^{-\sigma}\).
Put
\[
 Q=N^\beta,\qquad H=N^{\beta/2},\qquad
 \mathcal P_N=\{p\text{ prime}:P\le p\le Q\}.
\]
For \(A\subseteq[N]\), define
\[
 Y_l(N,A)=Z_{l,\mathcal P_N}
             \bigl(\mathcal L_{[N],\mathcal P_N,Q}\1_A\bigr).
\]
This is the quantity that will be compared between scales.

\begin{corollary}[Density domination and a uniform upper bound]
\label{lift:density}
For all positive integers \(l,N\) and all \(A\subseteq[N]\),
\[
 Y_l(N,A)\ge\left(\frac{|A|}{N}\right)^d\ge0.
\]
For every \(\epsilon>0\), there is a constant \(C_\epsilon\),
independent of \(l,N,A\), such that
\[
 Y_l(N,A)\le C_\epsilon N^\epsilon.
\]
\end{corollary}

\begin{proof}
The lift is real and has uniform mean \(|A|/N\).  The lower bound
follows from Lemma~\ref{lift:seminorm}.  Ordinary H\"older and the
uniform single-slot marginals give
\[
 Y_l(N,A)\le
 \left\|\mathcal L_{[N],\mathcal P_N,Q}\1_A\right\|_d^d.
\]
The integer \(d\) is even.  Also \(6\beta<1\), so for all sufficiently
large \(N\) one has \(N\ge10Q^6\).  Apply~\eqref{eq:12} with moment
order \(d\), choosing its exponent small enough that its \(d\)-th
power is at most \(C_\epsilon N^\epsilon\).  In the bounded remaining
range of \(N\), the lift has a bounded number of coefficients, each of
modulus at most one.  The probability normalization then supplies a
uniform bound independent of \(l\), and enlarging \(C_\epsilon\)
finishes the proof.
\end{proof}
